\section{Planar compactness and regularity facts}
\label{sec:known-facts}

We recall the precise form of the planar theory used below.  The distinction
between an absolute minimizer and a generalized limit matters only until we
know that the complement of the limiting set is connected.

For a closed rectifiable set $H\subset\mathbb R^2$ of locally finite length and a
function $v\in W^{1,2}(D\setminus H)$ for each bounded open $D$, define
\[
 E_0(v,H;V)=\int_{V\setminus H}|\nabla v|^2\,dx
                 +\mathcal H^1(H\cap V).
\]
An absolute minimizer of $E_0$ satisfies the comparison inequality for
all pairs in this local class agreeing outside a compact subset of $V$,
for every bounded open $V\Subset\mathbb R^2$.

In the compactness theorem, $(u_j,K_j)$ denotes a sequence of rescaled
minimizing pairs and $H$ is the local Hausdorff limit of their closed sets.
Write $D_k$, with $k$ in a countable index set, for the connected components of the limiting complement and choose
reference points $z_k\in D_k$.  Local Hausdorff convergence places each
fixed $z_k$ off $K_j$ for all sufficiently large $j$, where the regular
representative of $u_j$ is available.  On each $D_k$ the function $u_j$ is normalized
by subtracting $u_j(z_k)$.  In addition to the normalized
limiting harmonic functions, retain the limits
\[
 p_{kl}=\lim_{j\to\infty}\bigl(u_j(z_k)-u_j(z_l)\bigr)
 \in[-\infty,\infty].
\]
Write $v$ for the function given by these normalized limits on the
components.  The resulting object is a triple $(v,H,\{p_{kl}\})$.  Throughout this paper,
\emph{generalized minimizer} means a limit of \emph{reduced absolute} minimizers in the
sense of \cite[Definition~2.2.4]{delellis_focardi2025}; a \emph{global}
generalized minimizer has limiting domain $\mathbb R^2$.  We use this specific
class, including its information about finite relative constants.

\pagebreak[3]
\begin{proposition}[The planar theory used here]
\label{prop:planar-facts}
The following facts hold.
\begin{enumerate}
\item\label{item:blowup}
Let $(u,K)$ be a reduced, rectifiable absolute minimizer of the fidelity
functional in an open set $\Omega\subset\mathbb R^2$, with bounded fidelity datum and locally
bounded $u$.  For each $x\in K$ and each sequence $r_j\downarrow0$, a subsequence
of
\[
 K_j=\frac{K-x}{r_j},\qquad
 u_j(y)=\frac{u(x+r_jy)}{\sqrt{r_j}}
\]
on the expanding domains $(\Omega-x)/r_j$ has a global generalized limit
$(v,H,\{p_{kl}\})$.  The sets converge locally in
the bilateral Hausdorff sense.  The set $H$ is closed, rectifiable, locally of
finite length, and is the support of $\mathcal H^1\!\llcorner H$; moreover
$0\in H$.  For every bounded open $D\Subset\mathbb R^2$,
$v\in W^{1,2}(D\setminus H)$.

\item\label{item:connected-absolute}
If a global generalized minimizer has connected complement
$D=\mathbb R^2\setminus H$, then $(v,H)$ is an absolute minimizer of the
homogeneous energy $E_0$ under every compactly supported change.  In particular, changes whose support
meets $H$ are allowed.

\item\label{item:global-estimates}
Every global generalized minimizer has a reduced, closed, rectifiable set
$H$ of locally finite length, and $v\in W^{1,2}(D\setminus H)$ for every
bounded open $D\Subset\mathbb R^2$.  There are positive dimensional constants
$c,C$ such that
\[
 c r\leq\mathcal H^1(H\cap B_r(y))\leq C r
 \quad(y\in H,\ r>0),
\]
\[
 \int_{B_R\setminus H}|\nabla v|^2\,dx\leq 2\pi R
 \quad(R>0).
\]
Almost every point of $H$ with respect to $\mathcal H^1$ has a neighborhood in
which the entire set is a single $C^1$ arc.

\item\label{item:classification}
A global generalized minimizer has at most one unbounded connected component
of its closed set.  If its closed set disconnects the plane, that set is a
line or three half-lines meeting at $120$ degrees.  If all but at most one
connected component of its closed set lie in one compact set, the set is empty,
a line, three half-lines meeting at $120$ degrees, or a half-line.

\item\label{item:epsilon}
There are $\alpha\in(0,1)$ and $\varepsilon>0$ with the following property, with the constants
chosen for the fixed bounds on the fidelity datum.  Suppose that a reduced
absolute minimizer is defined in a neighborhood of $\overline{B_{2s}(x)}$,
where $s\leq1$.  If its closed set in $B_{2s}(x)$ has bilateral Hausdorff
distance less than $\varepsilon s$ from a diameter, a radius, or three radii
meeting at $120$ degrees, then its entire intersection with $B_s(x)$ is
$C^{1,\alpha}$-diffeomorphic to that model.  Any three arcs meeting in the
smaller ball meet at $120$ degrees.  The endpoint or junction in this conclusion
need not be the original center $x$.
\end{enumerate}
\end{proposition}

All references in this paragraph are to the 4 January 2025 author version
of De Lellis--Focardi \cite{delellis_focardi2025}.
Item~\ref{item:blowup} uses Assumption~2.2.1, Theorem~2.2.3 and
Definition~2.2.4, together with the lower density bound in Theorem~2.1.3.
Item~\ref{item:connected-absolute} specializes Definition~2.2.2(ii) and
Theorem~2.2.3(ii). For Item~\ref{item:global-estimates}, Lemma~2.1.2 and
equation~(1.4.1) give the total homogeneous energy bound $2\pi R$;
Theorem~2.1.3 gives lower density. The almost-everywhere arc assertion
follows from Theorem~1.5.2(i),(v), which applies to generalized minimizers:
the irregular part has dimension less than one, and the triple junctions
and regular loose ends form a discrete set. These sets have zero length;
every remaining point has a neighborhood consisting of a single regular arc.
For Item~\ref{item:classification}, Corollary~4.4.3 bounds the number of
unbounded components, Theorem~4.4.1 treats a disconnected complement,
and Theorem~1.4.6 treats a common compact remainder.
Item~\ref{item:epsilon} is Theorem~1.3.3, including its endpoint and
junction conclusions. We explain two applications whose hypotheses will
be important.

For \ref{item:blowup}, fix an ambient open set $W\Subset\Omega$ containing
$x$.  Apply the compactness theorem to the constant sequence of original
minimizers on $W$, with $\lambda_j=1$, $x_j=x$, and $r_j\downarrow0$.
Local boundedness supplies the uniform bound on the \emph{unscaled} functions
required by Assumption~2.2.1 of the reference; the rescaled domains exhaust
$\mathbb R^2$. Theorem~2.2.3 supplies $v\in W^{1,2}(D\setminus H)$
on each bounded open $D$, including square integrability up to $H$. This
stronger conclusion is part of the compactness theorem, not a consequence
of componentwise local convergence alone. It permits multiplication by smooth
cutoffs whose support meets $H$, as required in Section~\ref{sec:compact-piece}.
The density lower bound and length convergence ensure that the
Hausdorff limit remains reduced.  Indeed, at a limit point choose convergent
points of $K_j$, apply the lower bound in a small ball around them, and pass to
a slightly larger ball whose boundary has zero limiting length.  Thus every
neighborhood of that limit point has positive length.  The inclusion $0\in H$
also follows directly from $0\in K_j$.

For \ref{item:connected-absolute}, recall that a topological competitor in
\cite[Definition~2.2.2]{delellis_focardi2025} must preserve the separation
of exterior points which belong to different components of the original
complement. When that complement is connected, there are no such pairs
of points. Every admissible compact change therefore satisfies the
separation condition, and Theorem~2.2.3(ii) of that reference gives absolute
minimality on the whole plane. In particular, the support of a change may
meet $H$. This is the comparison class used in the compact-piece argument;
no condition of vanishing trace on $H$ is imposed.
